% ====================================================================
% FILE: content/13_branching_topology_de.tex
% Ontologische Verzweigung und globale Topologie
% Vollversion fuer Core 1.3 wiederhergestellt
% ====================================================================

\section{Ontologische Verzweigung und globale Topologie}
\label{sec:branching-topology}

Dieser Abschnitt rekonstruiert die Theorie der ontologischen Verzweigung und der
globalen Topologie von Kontinua. Er elaboriert Axiom~21 (Ontologische Verzweigung)
und integriert es mit der dimensionsbezogenen Hierarchie \(K_0\)–\(K_{10}\), den
Geburtsoperatoren \(\Psi_{x\to x+1}\) (Abschnitt~\ref{sec:operators-full}) und der
Kollaps--Wiedergeburt-Struktur (Abschnitt~\ref{sec:collapse-rebirth}).

Informell ist eine \emph{Verzweigung} ein maximaler Entwicklungspfad, auf dem ein
Kontinuum sukzessive dimensionale Geburten vollzieht, die mit einer gegebenen
Menge von Achsen und Einbettungsraeumen konsistent sind. Die Gesamtheit
solcher Verzweigungen besitzt eine natuerliche gerichtete Topologie, die eine
globale Sicht auf alle im Rahmen der OC-Axiome zulaessigen ontogenen Trajektorien
liefert.

\subsection{Definition ontologischer Verzweigungen}

Ein Kontinuumsniveau \(K_x\) ist durch
\[
  K_x =
  \big(
    \Omega_x, A_x, P_x(t), J_x(t),
    \Theta_x, \partial\Omega_x, C_x, k_x(t)
  \big)
\]
eingebettet in einen Raum \(M_x\) mit Achsenmenge \(A(M_x)\). Dimensionale Geburt
von \(K_x\) zu \(K_{x+1}\) wird vermittels eines Geburtsoperators
\(\Psi_{x\to x+1}\) (Abschnitt~\ref{sec:operators-full}).

\paragraph{Definition 13.1 (Ontologische Verzweigung).}
Eine \emph{ontologische Verzweigung} \(B\) ist eine endliche oder abzählbare
Sequenz von Kontinua
\[
  B = (K_{x_0}, K_{x_1}, K_{x_2}, \dots),
\]
zusammen mit Einbettungsraeumen \((M_{x_0}, M_{x_1}, \dots)\), so dass gilt:
\begin{enumerate}
  \item \textbf{Monotone Niveaus.}\quad
        \(x_0 < x_1 < x_2 < \dots\) und jeder Schritt
        \(K_{x_i} \to K_{x_{i+1}}\) ist ein dimensionaler Übergang
        (keine rein parametrische Evolution);
  \item \textbf{Geburt-Kompatibilitaet.}\quad
ae
        Für jedes \(i\) gibt es einen Geburtsoperator
        \(\Psi_{x_i\to x_{i+1}}\) mit
        \[
          (K_{x_{i+1}}, M_{x_{i+1}}) =
            \Psi_{x_i\to x_{i+1}}(K_{x_i}, M_{x_i}),
        \]
        wobei \(\Psi_{x_i\to x_{i+1}}\) die allgemeinen Geburtbedingungen aus
        Abschnitt~\ref{sec:operators-full} erfuellt;
  \item \textbf{Achsen-Monotonie.}\quad
        Die Achsenmengen wachsen monoton:
        \[
          A_{x_i} \subset A_{x_{i+1}}
          \subseteq A(M_{x_{i+1}}),
        \]
        und jede neue Achse ist nicht im strukturellen Spannungsraum der
        vorherigen Achsen darstellbar;
  \item \textbf{Einbettungs-Monotonie.}\quad
        Einbettungsraeume bilden eine Inklusionskette
        \[
          M_{x_0} \subset M_{x_1} \subset M_{x_2} \subset \dots
        \]
        mit \(A(M_{x_i}) \subseteq A(M_{x_{i+1}})\);
  \item \textbf{Maximalitaet.}\quad
        Die Verzweigung ist maximal mit Blick auf Erweiterung in steigender
        Niveau-Richtung: Es gibt kein Kontinuum \(K_{x^{\ast}}\) mit
        \(x^{\ast} > \sup_i x_i\), fuer das die Sequenz strikt erweitert werden
        kann unter Erhaltung der obigen Bedingungen.
\end{enumerate}

Wenn \(x_0 = 0\), ist die Verzweigung \emph{verwurzelt} im strukturellen
Substrat \(K_0\). Allgemeiner können Verzweigungen bezogen auf ein beliebiges
Ausgangsniveau \(K_{x_0}\) definiert werden, das als wirksame Wurzel fuer ein
restraines Bereichsgebiet dient (z. B. Chemie, Kognition, soziale Systeme).

\paragraph{Bemerkung 13.2 (Verzweigungs-Identitaet).}
Zwei Verzweigungen \(B\) und \(B'\) gelten als verschieden, wenn sie in
mindestens einem Punkt differieren:
\begin{itemize}
  \item in der Niveaereihe \(x_i\) (z. B. ausgelassene oder eingefuegte
        Ebenen),
  \item in den Achsenmengen \(A_{x_i}\) (verschiedene emergente Freiheitsgrade),
  \item in den Schwellenfamilien \(\Theta_{x_i}\) oder Einbettungsraeumen
        \(M_{x_i}\).
\end{itemize}
Verzweigungen kodieren daher \emph{Entwicklungsgeschichte(n)} im Raum der
Kontinua, nicht bloss statische Ebenensammlungen.

\subsection{Verzweigungstopologie}

Die Menge aller Verzweigungen ist mit einer natuerlichen gerichteten Struktur
versehen.

\paragraph{Definition 13.3 (Verzweigungsgraph).}
Sei \(\mathcal{K}\) die Klasse aller unter den OC-Axiomen zulaessigen
Kontinua und sei \(\mathcal{B}\) die Menge aller Verzweigungen. Der
\emph{Verzweigungsgraph} \(\mathcal{G}\) ist der gerichtete Graph, dessen:
\begin{itemize}
  \item Knoten Kontinua \(K_x\in\mathcal{K}\) sind,
  \item gerichtete Kanten \(K_x \to K_y\) den Geburtsoperatoren
        \(\Psi_{x\to y}\) mit \(y > x\) entsprechen.
\end{itemize}
Eine Verzweigung \(B\) ist dann ein gerichteter Pfad in \(\mathcal{G}\), der
bezueglich Erweiterung maximal ist.

Da Geburtsoperatoren die Dimensionalitaet nicht reduzieren koennen (Monotonie
der Dimension), ist \(\mathcal{G}\) ein gerichteter azyklischer Graph (DAG).
In der \emph{Zustandsdynamik} eines gegebenen Kontinuums \(K_x\) (uber interne
Zyklen \(C_x\)) koennen zwar Zyklen existieren, nicht jedoch in der
vertikalen Geburtstruktur.

\paragraph{Globale Topologie.}
Die globale topologische Struktur von \(\mathcal{G}\) ist eingeschraenkt durch:
\begin{itemize}
  \item \textbf{Gemeinsame Praefixe.}\quad
        Unterschiedliche Verzweigungen teilen gemeinsame Anfangssegmente, wenn sie
        in tieferen Ebenen uebereinstimmen (z. B. teilen alle physikalischen,
        chemischen und biologischen Verzweigungen \(K_0\)–\(K_2\));
  \item \textbf{Verzweigungs-Divergenz.}\quad
        Verzweigungen divergieren, wenn Geburtsoperatoren auf einem Niveau \(K_x\)
        verschiedene neue Achsen oder Schwellenstrukturen waehlen, was zu
distincten Weiterbildungen \(K_{x+1}^{(a)}\), \(K_{x+1}^{(b)}\), usw. führt;
  \item \textbf{Moegliche Fusionen.}\quad
        Durch Interaktionsoperatoren \(E_{\mathrm{int}}\) koennen unterschiedliche
        Verzweigungen spaeter Kontinua auf einer gemeinsamen Erweiterung erzeugen,
        wodurch eine DAG-artige statt rein baumartige Topologie entsteht;
  \item \textbf{Hypergraph-Struktur.}\quad
        Wenn Geburtsoperatoren mehrere Kontinua als Eingang erhalten (z. B.
        Fusionsereignisse), kann der Verzweigungsgraph zu einem gerichteten
        Hypergraphen erweitert werden mit Hyperkanten
        \((K_{x}^{(a)}, K_{x}^{(b)},\dots) \to K_{y}\).
\end{itemize}

Die Topologie der Verzweigungen interpoliert damit zwischen baumartigen
Strukturen (reine Spaltung, keine Fusion) und hypergraph-artigen Strukturen
(Verzweigungsfusion und kollektive Geburt). OC fixiert nicht eine einzige
globale Form, sondern beschraenkt alle zulaessigen Formen durch
Dimensions-Monotonie, Einbettungs-Kompatibilitaet und
Schwellenwertbedingungen.

\subsection{Verzweigungsbedingungen}

Verzweigung ist nicht an beliebiger Stelle moeglich; sie wird strukturell durch
spezifische Schwellen- und Einbettungsbedingungen ausgeloest.

\paragraph{Axiom 21 (Ontologische Verzweigung).}
Sei \(K_x\) ein Kontinuum auf Niveau \(x\) mit Einbettungsraum \(M_x\). Ein
Verzweigungsereignis in \(K_x\) ist erlaubt, wenn mindestens zwei unterschiedliche
Geburtsoperatoren
\[
  \Psi_{x\to x+1}^{(a)},\ \Psi_{x\to x+1}^{(b)},\dots
\]
existieren, so dass:
\begin{enumerate}
  \item \textbf{Unterscheidbare neue Achsen.}\quad
        Jeder Geburtsoperator waehlt eine strukturell verschiedene neue Achse
        \(A_{\mathrm{new}}^{(a)}\), \(A_{\mathrm{new}}^{(b)}\), \dots,
        die sich gegenseitig nicht im strukturellen Spannungsraum repräsentieren
        lassen, d. h.
        \[
          A_{\mathrm{new}}^{(a)} \notin \mathrm{span}_{\mathrm{str}}
            \big(A_x \cup \{A_{\mathrm{new}}^{(b)}\}\big), \ \text{etc.};
        \]
  \item \textbf{Dimensionale Spannung.}\quad
        Strukturelle Spannung ueberschreitet den dimensionalen Schwellenwert entlang
        mehrerer nicht-aequivalenter Unterschiedsrichtungen:
        \[
          T(K_x,t) > \Theta_{\mathrm{dim}}^{(a)}(K_x)
          \quad\text{and}
          T(K_x,t) > \Theta_{\mathrm{dim}}^{(b)}(K_x),
        \]
        wobei die Hochstellungen die unterschiedlichen dimensionalen Kanaele bezeichnen;
  \item \textbf{Einbettungs-Kapazitaet.}\quad
        Der Einbettungsraum \(M_x\) stuetzt alle Kandidatenachsen, d. h.
        \[
          A_{\mathrm{new}}^{(a)}, A_{\mathrm{new}}^{(b)}, \dots
          \in A(M_x);
        \]
  \item \textbf{Nichtleere zulaessige Regionen.}\quad
        Jedes Kandidatenkontinuum
        \(K_{x+1}^{(a)} = \Psi_{x\to x+1}^{(a)}(K_x, M_x)\), usw., besitzt ein
        nichtleeres zulaessiges Set \(\Omega_{x+1}^{(a)}\neq\emptyset\).
\end{enumerate}

In vielen konkreten Systemen werden diese Bedingungen realisiert durch:
\begin{itemize}
  \item Inkonsistenzen in Potenzialen \(P_x(t)\), die in einer einzelnen
        Erweiterung nicht vereinbar sind;
  \item Konkurrenz zwischen Zykluskomplexen \(C_x^{(a)}\), \(C_x^{(b)}\), deren
        gleichzeitige Stabilisierung die Stabilitaetsschwelle verletzen wuerde;
  \item Randgeometrien, die unterschiedliche, wechselseitig ausschliessende
        Wieder-Schliessungen zulassen (z. B. alternative Membranarchitekturen
        oder institutionelle Strukturen).
\end{itemize}

\paragraph{Schwellenwert der Verzweigung.}
Zur Bequemlichkeit kann ein \emph{Verzweigungs-Schwellenwert}
\(\Theta_{\mathrm{branch}}(K_x)\) als effektive Kombination aus dimensionalen und
Einbettungs-Schwellen eingefuehrt werden, so dass
\[
  T_{\mathrm{branch}}(K_x,t) > \Theta_{\mathrm{branch}}(K_x)
  \quad\Rightarrow\quad
  \text{Verzweigungsereignis moeglich.}
\]
Core~1.3 behandelt \(\Theta_{\mathrm{branch}}\) als abgeleitete Grösse statt als neue
primitive Schwellenklasse: Sie kodiert die gleichzeitige Erfuellung der Bedingungen
mehrerer Geburtsoperatoren.

\subsection{Verzweigungsinteraktion}

Verzweigungen entwickeln sich im Allgemeinen nicht isoliert. Wenn Kontinua aus
unterschiedlichen Verzweigungen in einem gemeinsamen Einbettungsraum koexistieren,
dann koennen sie durch den Interaktionsoperator \(E_{\mathrm{int}}\)
(Abschnitt~\ref{sec:operators-full}) wechselwirken.

\paragraph{Definition 13.4 (Verzweigungsinteraktion).}
Seien \(B^{(a)}\) und \(B^{(b)}\) zwei Verzweigungen mit Kontinua
\(\{K_x^{(a)}\}\) und \(\{K_y^{(b)}\}\), jeweils eingebettet in einen gemeinsamen
oder ueberlappenden Raum \(M\). Eine \emph{Verzweigungsinteraktion} ist eine
Familie von Interaktionsoperatoren
\[
\begin{aligned}
  E_{\mathrm{int}}^{x,y}:\quad
    &(K_x^{(a)}, K_y^{(b)}, M)\
    &\longrightarrow (K_x^{(a)\prime}, K_y^{(b)\prime}, M'),
\end{aligned}
\]
definiert fuer relevante Ebenenpaare \((x,y)\).

Typische Regime sind:

\begin{itemize}
  \item \textbf{Koexistenz.}\quad
        Fluesse sind schwach gekoppelt und erlauben beiden Verzweigungen, ihre
        Continuumness zu erhalten oder zu erhoehen:
        \[
          k^{(a)}(t+dt) \ge k^{(a)}(t),\quad
          k^{(b)}(t+dt) \ge k^{(b)}(t).
        \]
  \item \textbf{Interferenz.}\quad
        Querverzweigungs-Fluesse \(J^{(a\leftrightarrow b)}\) erzeugen
        strukturelle Spannung, die eine oder beide Verzweigungen naeher an die
        Todes-Schwellen dringt, ohne sofortigen Kollaps.
  \item \textbf{Suppression.}\quad
        Eine Verzweigung treibt die andere unter ihre Viabilitaetsschwellen,
        was zum Kollaps der unterdrueckten Verzweigung fuehrt, waehrend der
        Unterdruecker lebensfaehig bleibt oder sogar staerker wird.
  \item \textbf{Parasitismus.}\quad
a
        Eine Verzweigung entnimmt unterstuetzende Fluesse von einer anderen und
        steigert ihre eigene Continuumness, waehrend die Zykluskomplexe der
        anderen ausgehöhlt werden.
  \item \textbf{Fusion.}\quad
        Durch starke Kopplung und gemeinsame Randstrukturen koennen zwei
        Verzweigungen ein neues Kontinuum erzeugen, das auf einer gemeinsamen
        Erweiterungs-Verzweigung liegt. Formal entspricht dies einem
        zusammengesetzten Geburtsoperator, gespeist aus sowohl \(B^{(a)}\) als auch
        \(B^{(b)}\).
\end{itemize}

Gemeinsame Randstrukturen sind besonders wichtig. Wenn
\(\partial\Omega^{(a)}\cap\partial\Omega^{(b)}\neq\emptyset\), dann:
\begin{itemize}
  \item Querverzweigungs-Fluesse, begrenzt auf die Schnittmenge, koennen
        Kopplung vermitteln;
  \item Teilzustände auf dem gemeinsamen Rand koennen beide Verzweigungen
        synchronisieren oder destabilisieren;
  \item neue Kontinua koennen am gemeinsamen Rand entstehen und neue
        Verzweigungen oder Verzweigungsfusionszustaende initiieren.
\end{itemize}

\subsection{Kollaps von Verzweigungen}

Verzweigungen koennen in mehreren Varianten enden; so wird der Kollapsbegriff von
individuellen Kontinua (Abschnitt~\ref{sec:collapse-rebirth}) auf ganze
Entwicklungspfade erweitert.

\paragraph{Definition 13.5 (Verzweigungs-Kollaps).}
Eine Verzweigung \(B = (K_{x_0}, K_{x_1}, \dots)\) \emph{kollabiert} auf Niveau
\(x_n\), wenn:
\begin{enumerate}
  \item das Kontinuum \(K_{x_n}\) im Sinn von Definition~12.1 kollabiert
        (leere zulaessige Menge, \(k\to 0\), Schwellenverletzung);
  \item es kein Kontinuum \(K_{x_{n+1}}\) mit \(x_{n+1}>x_n\) gibt, auf
        dem ein Geburtsoperator \(\Psi_{x_n\to x_{n+1}}\) auf irgendeiner Restmenge
        von \(K_{x_n}\) definiert werden kann
        (Abschnitt~\ref{sec:collapse-rebirth});
  \item die Verzweigung nicht weiter erweitert werden kann, wobei die Axiome der
        dimensionalen Monotonie und der Einbettungs-Kompatibilitaet eingehalten
        bleiben.
\end{enumerate}

Drei strukturell unterschiedliche Szenarien sind relevant:

\begin{itemize}
  \item \textbf{Terminaler Kollaps.}\quad
        Das letzte Kontinuum auf der Verzweigung kollabiert und kein
        Wiederaufbau auf höheren Niveaus ist moeglich. Die Verzweigung ist
        endlich und endet bei \(x_n\).
  \item \textbf{Fusions-Kollaps.}\quad
        Die Verzweigung wird über Fusion in eine stabilere Verzweigung absorbiert:
        Kontinua in \(B\) verlieren ihre Identitaet als unabhängige Knoten im
        Verzweigungsgraph und werden Teilstruktur eines Kontinuums in einer
        anderen Verzweigung. Aus der Perspektive von \(B\) ist dies ein
        Kollaps durch Verlust der Verzweigungs-Identitaet.
  \item \textbf{Asymptotisches Ausdünnen.}\quad
        Die Verzweigung endet nicht in einem einzelnen dramatischen Kollaps,
        sondern geht in ein Regime ueber, in dem die höheren Niveaus zunehmend
        marginal sind (\(k_{x_i}\to 0\) fuer \(i\to\infty\)) und keine
        weiteren dimensionalen Geburten stützen. Die Verzweigung ist im Index
        unendlich, aber in effektiver Komplexitaet endlich.
\end{itemize}

\paragraph{Globale Restriktionen.}
Die OC-Axiome erzwingen Restriktionen, die unkontrollierte unendliche
Verzweigungs-Kaskaden auf endlichen Niveaus verhindern:
\begin{itemize}
  \item der Einbettungsraum auf jedem Niveau hat eine endliche Achsenkapazität
        je nach Domäne, wodurch die Zahl unterschiedlicher Geburtskanäle
        eingeschraenkt ist;
  \item dimensionsbezogene Schwellen erfordern nichttriviale strukturelle
        Spannung; wiederholte Verzweigung ohne ausreichend starke stützende
        Fluesse wuerde Verzweigungen in Richtung Kollaps bewegen;
  \item Monotonie-Saetze der Komplexitaet beschraenken viable Verzweigungen auf
        solche, deren Komplexitaetswachstum zu Verfuegung stehenden Zyklen und
        Schwellen kompatibel ist.
\end{itemize}

\subsection{Konsequenzen fuer die K-Level-Hierarchie}

Die Verzweigungsstruktur verfeinert die vertikale Hierarchie \(K_0\)–\(K_{10}\),
indem sie festlegt, welche Sequenzen von Ebenen und Achsen wirklich
realisierbar sind.

\paragraph{Ontogenetische Entwicklungspfade.}
Jede Verzweigung definiert einen \emph{ontogenetischen Entwicklungspfad}:
\begin{itemize}
  \item in physikalischen Domänen beschreiben Verzweigungen Wege, in denen feldtheoretische
        Kontinua auf \(K_2\) chemische, protocellulaere und biologische
        Kontinua auf \(K_3\)–\(K_5\) hervorbringen;
  \item in kognitiven und sozialen Domänen beschreiben Verzweigungen wie neuronale
        Kontinua \(K_5\) kognitive Kontinua \(K_6\) unterstuetzen koennen,
        die wiederum soziale und zivilisatorische Kontinua auf \(K_7\)–\(K_8\)
        stützen;
  \item in theoretischen Domänen verbinden Verzweigungen empirische Kontinua mit
        theoretischen und meta-theoretischen Kontinua \(K_9\)–\(K_{10}\).
\end{itemize}
Die Vielfalt beobachteter physikalischer, chemischer, biologischer und sozialer
Systeme entspricht der Vielzahl von Verzweigungen, die mit einem gemeinsamen
Substrat vereinbar sind.

\paragraph{Zulaessige vs. verbotene Übergänge.}
Die Verzweigungstopologie kodifiziert, welche Übergange zwischen Niveaus
erlaubt sind:
\begin{itemize}
  \item Kanten im Verzweigungsgraphen \(\mathcal{G}\) entsprechen zulaessigen
dimensionalen Geburtsoperatoren;
  \item das Fehlen einer Kante \(K_x \to K_y\) (mit \(y>x\)) drückt eine
        strukturelle Unmoeglichkeit aus: Kein Einbettungsraum und keine
        Schwellenlandschaft kann diesen direkten Übergang tragen, sofern nicht
        die erforderlichen Zwischenebenen durchlaufen werden;
  \item verbotene Übergänge verletzen typischerweise entweder die dimensionsbezogene
        Monotonie (versuchte Vereinfachung von Achsen) oder Einbettungs-Schwellen
        (versuchte Realisierung von Achsen ohne Stuetze in \(M_x\)).
\end{itemize}

\paragraph{Vertikale Kohärenz.}
Da Verzweigungen auf tiefere Niveaus reduzierbar sein muessen, wenn
zusätzliche Achsen eingefroren sind (Abschnitt~\ref{sec:operators-full}), erzwingt
die Verzweigungsstruktur vertikale Kohärenz:
\begin{itemize}
  \item jedes hoehere Kontinuum \(K_y\) reduziert zu einem niederen Kontinuum
        \(K_x\) entlang einer Verzweigung, wenn zusätzliche Achsen konstant
        gehalten werden;
  \item umgekehrt wirken untere Niveaus als notwendige Schichten fuer hoehere;
  \item Verzweigungen dokumentieren diese Kohärenz als gemeinsame Praefixe
        ueber Domänen (z. B. muessen alle biologischen Verzweigungen
        physikalische und chemische Ebenen enthalten).
\end{itemize}

\paragraph{Globales Bild.}
Zusammen mit der verticalen Hierarchie \(K_0\)–\(K_{10}\), der
Schwellenlandschaft und der Verzweigungstopologie bildet sich das globale
strukturelle Bild von OC.
Die operative Struktur wird von \(F\), \(G\), \(H\), \(Q\), \(R\),\
\(S\), \(U\), \(\Psi\) und \(E_{\mathrm{int}}\) getragen:
\begin{itemize}
  \item Kontinua sind lokale Objekte mit interner Dynamik;
  \item Verzweigungen sind globale Entwicklungspfade durch die Ebenen;
  \item der Verzweigungsgraph kodiert, welche Pfade strukturell erlaubt sind;
  \item Kollaps und Wiedergeburt schneiden und rekonfigurieren Verzweigungen
        ueber die Zeit.
\end{itemize}

Core~1.3 fixiert dieses globale Bild auf strukturellem Niveau. Spätere Erweiterungen
und domänenspezifische Studien können Verzweidigungswahrscheinlichkeiten
quantifizieren, konkrete Beispiele von Verzweigung in konkreten Systemen explizit
untersuchen und empirische Kriterien zur Identifikation von Verzweigungsstruktur
in Realwelt-Daten entwickeln.

% END OF FILE

